\begin{claim}{F.24}{Proved}{All-field rank-one cocycles and the actual marginal affine fixed points}
Retain the actual continuous inclusions
$\mathcal K_\pm\leq R_\pm\leq V_\pm$ and the source maps of
F.19--F.23. Let $k$ be a finite field, with its discrete topology,
and let $\chi:V_\pm\to k^\times$ be a continuous character whose
restriction to $R_\pm$ is nontrivial. Give the additive group of
$k$ the left action $g\cdot b=\chi(g)b$, and denote this module
by $k_\chi$. The actual restriction maps are injective:
\[
\begin{aligned}
 Z^1_{\mathrm{cont}}(V_\pm,k_\chi)&\longrightarrow
 Z^1_{\mathrm{cont}}(\mathcal K_\pm,k_{\chi|\mathcal K_\pm}),\\
 H^1_{\mathrm{cont}}(V_\pm,k_\chi)&\longrightarrow
 H^1_{\mathrm{cont}}(\mathcal K_\pm,k_{\chi|\mathcal K_\pm}).
\end{aligned}                                                   \tag{F.24.1}
Here a crossed homomorphism satisfies
$d(gh)=d(g)+\chi(g)d(h)$, and the principal crossed
homomorphisms are $d_b(g)=(\chi(g)-1)b$, for $b\in k$.

Write $\operatorname{Aff}_1(k)=k\rtimes k^\times$, acting on
$k$ by $(b,a)z=b+az$, and let
$T_k=\{(b,1):b\in k\}$ be its translation subgroup.
For every continuous homomorphism
$\rho:V_\pm\to\operatorname{Aff}_1(k)$ one has
\[
 \operatorname{Fix}_k\rho(\mathcal K_\pm)
       =\operatorname{Fix}_k\rho(R_\pm).                        \tag{F.24.2}
\]
For the representations whose scalar character is nontrivial
on $R_\pm$, both sides also equal
$\operatorname{Fix}_k\rho(V_\pm)$.
For every representation in the specified class
\[
 \left\{\rho\in\operatorname{Hom}_{\mathrm{cont}}
       (V_\pm,\operatorname{Aff}_1(k)):
           \rho(\mathcal K_\pm)\cap T_k=\{1\}\right\},
\]
the exact marginal images agree:
\[
                 \rho(\mathcal K_\pm)=\rho(R_\pm).             \tag{F.24.3}
\]
The coefficient calculation proving (F.24.1) has integer
Laurent coefficients and proves its rank assertion over every
field, including characteristics two and three. Its input is
the literal negative marginal braid data of F.23, with the
same source lassos and comparison maps. The positive assertion
is transported by the explicit F.23.13 isomorphism.
The full marginal equalities and the equality $K^+=D^+$
remain unresolved.
\end{claim}

\begin{proof}
\emph{The retained words and the exact coefficient maps.}
We first use the negative coordinates
\[
 \begin{gathered}
 V_-=\widehat F(X_1,X_2,X_3,X_4,X_5,X_6,X_7),\qquad
 X_{j+3}=U_j\quad(1\leq j\leq4),\\
 R_-=\ker\bigl(V_-\longrightarrow
                     \widehat F(X_1,X_2,X_3)\bigr).
 \end{gathered}                                                \tag{F.24.4}
These are the actual F.22.5 filling coordinates. Put
$\Lambda=\mathbb Z[t_1^{\pm1},\ldots,t_7^{\pm1}]$.
Let $F_7$ be the discrete free group on these same seven
generators, with its dense completion map $F_7\to V_-$.
The universal character and crossed homomorphism are
\[
 t:F_7\longrightarrow\Lambda^\times,\quad t(X_i)=t_i,
 \qquad
 D:F_7\longrightarrow\Lambda^7,\quad D(X_i)=e_i,
                                                               \tag{F.24.5}
\]
where $\Lambda^7$ consists of row vectors and
$D(vw)=D(v)+t(v)D(w)$. These maps are well defined on
reduced words: the inverse formula is
$D(v^{-1})=-t(v)^{-1}D(v)$, which cancels an adjacent
letter and its inverse. Equivalently they are the two
coordinates of the group homomorphism
$F_7\to\Lambda^7\rtimes\langle t_1,\ldots,t_7\rangle$
specified on its free generators.

For a character $\chi$ and a continuous crossed homomorphism
$d$ as in the statement, let
\[
 \operatorname{ev}_\chi:\Lambda\longrightarrow k,
       \qquad t_i\longmapsto\chi(X_i),\qquad
 v_d=(d(X_1),\ldots,d(X_7))^{\mathsf T}.
\]
This is a ring homomorphism because each $\chi(X_i)$ is
nonzero. Induction on the reduced word gives the exact formula
\[
             d(w)=\operatorname{ev}_\chi(D(w))v_d
                      \qquad(w\in F_7).                       \tag{F.24.6}
\]
All words below are finite words in this dense free group,
interpreted in $V_-$ by its stated completion map.
In particular, the calculation does not substitute a
completion of a different source subgroup.

Let $\nu_0,\ldots,\nu_7$ be the eight negative collision
words in F.23, in the order
\[
 (u,\epsilon)=
 (1,+),(1,-),(\zeta^2,+),(\zeta^2,-),
 (\zeta^3,+),(\zeta^3,-),(\zeta^4,+),(\zeta^4,-).
                                                               \tag{F.24.7}
\]
The sign here selects the two roots of each negative
collision square $B_u$; every word in (F.24.7) belongs
to the negative marginal. The original paths $\eta$,
$G$, $T$, and $\Xi$ and their affine normalization are
exactly those retained in F.23.19--22.
Let $B_i:V_-\to V_-$, $0\leq i\leq10$, be conjugation
on this point-pushing group by the image of the $i$th
old source generator under the geometric source section
$j_{\mathrm{geo}}$ of F.23.23. Thus $B_i$ is a continuous
automorphism and
\[
 B_i(\mathcal K_-)=\mathcal K_- .                              \tag{F.24.8}
\]
Indeed the source filling kernel is normal in the source;
its image is $K^+$; projection to the negative marginal
commutes with this conjugation. These are the original
maps, with the old and geometric section comparison
$j_{\mathrm{geo}}(l)=c(l)j(l)$ from F.23.23--24.
The old source indices are $0$ for its zero puncture,
$1,\ldots,5$ for $+\zeta^j$, and $6,\ldots,10$ for
$-\zeta^j$, in each case $0\leq j\leq4$.

Every $B_i$ preserves the abelianization of $V_-$.
For the explicit proof in these coordinates, start the
decorated Artin recurrence F.23.26--28 with the first
seven characters $t_1,\ldots,t_7$ and the eighth character
one. A crossing exchanges the two characters, and each
recorded old braid is pure. Its final conjugation by $u$
does not change them. Hence $t(B_i(X_j))=t_j$, and
\[
 A_i=\bigl(D(B_i(X_j))\bigr)_{1\leq j\leq7},\qquad
 D(B_i(w))=D(w)A_i .                                          \tag{F.24.9}
\]
Both identities follow on generators and then on products
and inverses from the crossed-homomorphism identity.
Consequently the following sixty-four rows evaluate
the crossed homomorphism on actual elements of
$\mathcal K_-$:
\[
 \begin{gathered}
 a_r=D(\nu_r)\quad(0\leq r\leq7),\qquad
 a_rA_i\quad(0\leq r\leq7,\quad i\in\mathscr I),\\
                 \mathscr I=(2,9,8,10,1,0,7).
 \end{gathered}                                                \tag{F.24.10}
\]
Number the rows from zero, first the eight $a_r$, then
the eight rows for each index of $\mathscr I$ in the
displayed order.

\emph{Exact production and rank of the sixty-four rows.}
For completeness we give the full arithmetic recipe
and the terminal identities of its certificate. Encode
an element of the universal affine group by $(a,v)$,
where $a\in\mathbb Z^7$ is its character exponent and
$v\in\Lambda^7$ its crossed-homomorphism row. Then
\[
 (a,v)(b,w)=(a+b,v+t^a w),\qquad
 (a,v)^{-1}=(-a,-t^{-a}v).                                    \tag{F.24.11}
\]
Initialize the eight Artin images as $(e_j,e_j)$ for
$1\leq j\leq7$ and $(0,0)$ for the erased eighth
generator, and initialize all eight conjugators as
$(0,0)$. Apply the exact positive and inverse substitutions
F.23.26--27 in this group, using (F.24.11). Put
$u=c_8^{-1}$ at the end. The row of $u$ is the seed row;
the rows of $u x_j u^{-1}$ form $A_i$ for an old word.
This applies the proved decorated recurrence to a
specified group, so it evaluates the same actual
words as F.23.28. The separate producer records
eleven old matrices and eight seed rows. All of them
also satisfy
\[
 D(w)(t_1-1,\ldots,t_7-1)^{\mathsf T}=t(w)-1,                 \tag{F.24.12}
\]
which follows directly from (F.24.5).

Here is the exact row elimination used to check their
rank. In chart $j$, work over
$\Lambda[(t_j-1)^{-1}]$. A row normalization multiplies
by a Laurent monomial and divides by a common power of
$t_j-1$. A pivot $p$ is accepted only with the exact
coefficient identity
$p=\varepsilon t^a(t_j-1)^m$, where
$\varepsilon\in\{1,-1\}$ and $m\geq0$. With coefficient
$c$ in its column, replace a later row $r$ by $pr-c r_p$
and normalize again. Each operation has an invertible
row multiplier in the displayed ring and preserves
the row module there. Every polynomial is stored as
its integer coefficients and seven integer exponents;
the record retains the input coefficients, shifts,
division counts, pivot columns, and resulting rows.

The following are the complete pivot orders, with
zero-based row indices and one-based column indices.
The last chart first specializes $t_4=t_5=t_6=1$.
Each listed final pivot coefficient is a signed Laurent
monomial, so every pivot remains nonzero in every field
on its chart.
\[
\begin{array}{c|l|l|r}
 j&\text{pivot rows}&\text{pivot columns}&\text{row operations}\\\hline
 4&11,8,0,6,4,1&5,4,1,7,6,3&389\\
 5&11,8,3,6,4,2,34&5,4,1,7,6,2,3&444\\
 6&11,8,13,14,20,5,28&5,4,6,7,2,1,3&445\\
 7&0,2,4,14,39,6,46&4,5,6,7,2,1,3&196
\end{array}                                                     \tag{F.24.13}
\]
The first chart has one remaining column, column two.
After its six displayed pivots, the retained rows
numbered five and forty-eight are exactly
\[
 \begin{aligned}
 r_5&=(1-t_6)e_2,\\
 r_{48}&=\bigl(t_1t_2t_3t_4t_5t_7(t_6-1)-1\bigr)e_2.
 \end{aligned}
\]
They give the explicit seventh pivot
\[
           r_{48}+t_1t_2t_3t_4t_5t_7 r_5=-e_2 .              \tag{F.24.14}
\]
Thus all four charts have rank seven. The first three
charts cover $t_4\ne1$, $t_5\ne1$, and $t_6\ne1$.
The remaining stratum has those three values equal to
one and $t_7\ne1$, and is covered by the fourth chart.
This proves rank seven at every seven-tuple of nonzero
field elements for which some $t_j\ne1$, $4\leq j\leq7$.
No integer other than $1$ or $-1$ is inverted.

The accompanying exact certificate contains every row
and operation just described. Its independent verifier
imports none of the three producers: it reconstructs
the nineteen word evaluations with the group law
(F.24.11), reconstructs all sixty-four rows, replays
every row operation, proves every division by multiplying
back by $t_j-1$, and checks the triangular pivot orders
and (F.24.14). It also replays the longer scalar-ideal
certificate retained from the original search. There
are $13\,982$ exact assertions in this independent
verification. The manifest fixes the braid input,
all producers and outputs, and this verification by
size and SHA-256. These finite polynomial identities,
together with their invertible operations, establish
the asserted rank over all characteristics.

\emph{Restriction, continuity, and the sign comparison.}
The restriction $\chi|R_-$ is nontrivial precisely when
at least one of $\chi(X_4),\ldots,\chi(X_7)$ differs
from one: $R_-$ is the closed normal subgroup generated
by these four generators, and the target $k^\times$
is abelian. For a crossed homomorphism vanishing on
$\mathcal K_-$, every row in (F.24.10) annihilates $v_d$
by (F.24.6) and (F.24.8)--(F.24.9). The rank assertion
therefore gives $v_d=0$. The crossed-homomorphism rule
then gives zero on the dense group $F_7$; continuity
into the finite Hausdorff space $k$ gives zero on
$V_-$. This proves injectivity on $Z^1$.
For a class restricting to zero in $H^1$, its
representative restricts to $(\chi-1)b$ for some $b\in k$.
Subtract that same principal crossed homomorphism on
$V_-$ and apply the proved $Z^1$ injectivity. This
proves injectivity on $H^1$ with the exact coefficient
restriction in (F.24.1).

For the positive sign use the specific continuous
isomorphism $\mathcal B_+:V_-\to V_+$ of F.23.13.
It carries $R_-$ onto $R_+$ and $\mathcal K_-$ onto
$\mathcal K_+$. Pullback sends $(\chi,d)$ to
$(\chi\mathcal B_+,d\mathcal B_+)$; pullback along
$\mathcal B_+^{-1}$ is its inverse. It preserves the
crossed-homomorphism equations and the principal
rows and commutes literally with restriction.
Nontriviality on $R_+$ transfers to nontriviality on
$R_-$. These maps prove both positive assertions.

\emph{The affine representations and their fixed points.}
Write a given representation as
$\rho(g)=(d(g),\chi(g))$. Its multiplication law makes
$d$ a continuous crossed homomorphism. For a point
$z\in k$ put
\[
       d_z(g)=\rho(g)z-z=d(g)+(\chi(g)-1)z .                  \tag{F.24.15}
\]
It is another continuous crossed homomorphism for
the same character. A point fixed by
$\rho(\mathcal K_\pm)$ makes $d_z$ vanish there.
For $\chi|R_\pm$ nontrivial, (F.24.1) makes $d_z$
zero on $V_\pm$. The inclusions of fixed sets from
$\mathcal K_\pm\leq R_\pm\leq V_\pm$ now give the
three equal fixed sets claimed above.

For $\chi|R_\pm=1$, the image of $R_\pm$ lies in
the abelian group $T_k$. Apply the F.21 normal-generation
conclusion through the actual surjection
$p_\pm:D^+\to R_\pm$ of F.22.6 and then through
$\rho|R_\pm$. The normal closure of
$\rho(\mathcal K_\pm)$ in $\rho(R_\pm)$ is
$\rho(R_\pm)$; in this abelian group the closure is
the subgroup itself. Hence their images agree and
their fixed sets agree. This completes (F.24.2).

Finally put $M=\rho(\mathcal K_\pm)$ for a representation
in the class stated before (F.24.3). Its scalar
projection $M\to k^\times$ is injective because its
kernel is $M\cap T_k$. Thus $|M|$ divides $|k|-1$
and is invertible in $k$. The explicit point
\[
                     z_M=|M|^{-1}\sum_{m\in M}m(0)            \tag{F.24.16}
\]
is fixed by $M$. For $h=(b,a)\in M$, applying its
affine formula to this average gives
$h(z_M)=|M|^{-1}\sum_{m\in M}(hm)(0)=z_M$, because
left multiplication by $h$ permutes $M$.
For a nontrivial character on $R_\pm$, (F.24.2) and
its stronger assertion put all of $\rho(V_\pm)$ in
the stabilizer of $z_M$. The scalar map on that
stabilizer has the explicit inverse
$a\mapsto((1-a)z_M,a)$, so the stabilizer is abelian.
F.21 normal generation, transported through the same
two actual maps as above, gives (F.24.3). For a
trivial character on $R_\pm$, the preceding abelian
translation argument already gives (F.24.3).

The remaining finite-affine case consists exactly of
representations with $\rho(\mathcal K_\pm)\cap T_k$
nontrivial. The argument has not proved equality of
the two marginal images in that class. In particular
it has not constructed an ambient invariant quotient
that kills this translation intersection. The actual
inclusions of F.23.1 and the interaction kernel of
F.22.6--8 remain the maps whose surjectivity is under
investigation.
\end{proof}
